\section{Generality and consequences}\label{sec:consequences}

The construction and monomial extraction prove Artinian EGH over
\(\C\). We first extend that Hilbert-function conclusion to the scope
of Corollary~\ref{cor:characteristic-zero}, and then derive two further
consequences. These arguments use the EGH conclusion rather than any
additional division-algebra construction.

\subsection{Regular sequences over characteristic-zero fields}
\label{subsec:generality}

\begin{proof}[Proof of Corollary~\ref{cor:characteristic-zero}]
If \(I=S_F\), take \(J=S_F\). Otherwise, first work over \(\C\).
Choose a linear coordinate change for which
\(f_1,\ldots,f_c,x_{c+1},\ldots,x_n\) is regular. The images
\(\overline f_1,\ldots,\overline f_c\) in
\(\C[x_1,\ldots,x_c]\) form a full regular sequence.
Corollary~\ref{thm:main} applies to these images and, more generally,
to every full regular sequence of the prescribed degrees in \(c\)
variables and every homogeneous ideal containing it. This universal
assertion is the premise of the Artinian reduction of
Caviglia--Maclagan \cite[Proposition~10]{CM08}, which gives the desired
Hilbert-function comparison for every ideal containing the original
sequence in \(n\) variables. The lexicographic embedding modulo \(P_F\) gives the stated
lex-plus-powers choice; this reduction and its lex formulation are also
recalled in \cite[Section~4]{CK14}.

For general \(F\), choose homogeneous generators of \(I\) and let
\(F_0\subseteq F\) be the subfield generated over \(\Q\) by their
coefficients and those of the \(f_i\). Let \(I_0\subseteq F_0[x]\) be
generated by those generators together with the \(f_i\). Then
\(I_0F[x]=I\), and faithful flatness of \(F/F_0\) descends regularity
of the sequence to \(F_0[x]\). Embed the finitely generated field
\(F_0\) into \(\C\) and apply the complex case after scalar extension.
Define a monomial ideal over \(F_0\), and then over \(F\), using the same
monomial exponents as the resulting complex ideal. Dimensions of graded
pieces commute with both field extensions from \(F_0\), so the Hilbert
equalities descend to \(F\). The same observation identifies the
degreewise lex-plus-powers choice.
\end{proof}

The companion paper proves the stronger graded Betti inequalities in
this same characteristic-zero scope
\cite[Corollary~1.2]{BettiCompanion}, using the Artinian LPP reduction
of Caviglia--Kummini \cite[Theorem~4.1]{CK14}. The construction in
Theorem~\ref{thm:construction} itself remains the stated construction
for a full complex regular sequence.

\subsection{Local cohomology}

Caviglia--Sbarra's conditional theorem
\cite[Theorem~4.4]{CavigliaSbarra13} also gives a local-cohomology
comparison for the same lex-plus-powers ideal.

\begin{corollary}[Lex-plus-powers extremality for local cohomology]
\label{cor:local-cohomology}
Let \(F\) be a field of characteristic zero, let
\(S_F=F[x_1,\ldots,x_n]\) be standard graded, and put
\(\mathfrak m=(x_1,\ldots,x_n)\). Let
\[
 1\le c\le n,\qquad 2\le a_1\le\cdots\le a_c,
 \qquad P_F=(x_1^{a_1},\ldots,x_c^{a_c}).
\]
Suppose that a homogeneous ideal \(I\subseteq S_F\) contains a homogeneous
regular sequence \(f_1,\ldots,f_c\) with \(\deg f_k=a_k\). Let \(J\) be the
lex-plus-powers ideal prescribed by Corollary~\ref{cor:characteristic-zero}
for \(P_F\), the order \(x_1>\cdots>x_n\), and the Hilbert function of
\(S_F/I\). For this single ideal,
\[
 \dim_F H^i_{\mathfrak m}(S_F/I)_j
 \le
 \dim_F H^i_{\mathfrak m}(S_F/J)_j
 \qquad(i\ge0,\ j\in\Z).
\]
\end{corollary}

\begin{proof}
If \(I=S_F\), both sides vanish. Otherwise, since \(F\) is infinite,
complete \(f_1,\ldots,f_c\) to a regular sequence by \(n-c\) linear forms.
After a linear change of coordinates, the images
\(\overline f_1,\ldots,\overline f_c\) in
\(\overline S=F[x_1,\ldots,x_c]\) form a full regular sequence of the same
degrees. The list of added linear forms is empty when \(c=n\).

Corollary~\ref{cor:characteristic-zero}, applied in \(\overline S\) to every
homogeneous ideal containing this sequence, supplies precisely the Artinian
EGH premise of Caviglia--Sbarra \cite[Theorem~4.4]{CavigliaSbarra13}.
Their Clements--Lindstr\"om image of \(I/(f_1,\ldots,f_c)\) in \(S_F/P_F\)
has preimage \(J\) in \(S_F\). The base-independence identification used
there gives local cohomology with support in \(\mathfrak m\), and their
Hilbert-series comparison is coefficientwise in the full internal
\(\Z\)-grading for every cohomological degree \(i\ge0\). This gives the
displayed inequalities directly over \(F\).
\end{proof}

When \(c<n\), this statement includes positive-dimensional quotients.
When \(c=n\), both quotients are Artinian: their zeroth local cohomology
is the quotient itself and their higher local cohomology vanishes, so the
comparison reduces to the Hilbert-function equality.

\subsection{Quadratic Cayley--Bacharach}

The quadratic case also gives the following scheme-theoretic
Cayley--Bacharach consequence.

\begin{corollary}[Quadratic Cayley--Bacharach]
\label{cor:quadratic-cayley-bacharach}
Let \(F\) have characteristic zero, let \(r\ge1\), and let
\(\Omega\subseteq\PP^r_F\) be a zero-dimensional scheme-theoretic complete
intersection of \(r\) quadrics, of degree \(2^r\). For an integer
\(1\le k\le r\) and a degree-\(k\) hypersurface \(X\subseteq\PP^r_F\),
\[
 \Omega\not\subseteq X
 \quad\Longrightarrow\quad
 \deg(\Omega\cap X)\le 2^r-2^{r-k}.
\]
Equivalently, if \(X\) contains a closed subscheme of \(\Omega\) of degree
strictly greater than \(2^r-2^{r-k}\), then it contains \(\Omega\)
scheme-theoretically. Degrees here mean scheme lengths over \(F\);
no reducedness assumption is needed.
\end{corollary}

\begin{proof}
First let \(m\ge0\) and let \(\Gamma\subseteq\Omega\) fail to impose
independent conditions on degree-\(m\) forms, so that
\(H_\Gamma(m)<\deg\Gamma\), where
\(H_\Gamma(d)=\dim_F(S/I_\Gamma)_d\),
\(S=F[x_0,\ldots,x_r]\), and \(I_\Gamma\) is the saturated homogeneous
ideal. Choose a linear form \(\ell\) nonvanishing on \(\Omega\). It is a
nonzerodivisor on the homogeneous coordinate rings of both \(\Omega\)
and \(\Gamma\). Thus the images of the defining quadrics form a regular
sequence in \(S/(\ell)\cong F[x_1,\ldots,x_r]\), and, for
\(A=S/(I_\Gamma,\ell)\),
\[
 \deg\Gamma=\sum_{j\ge0}\dim_F A_j,
 \qquad H_\Gamma(m)=\sum_{j=0}^m\dim_F A_j.
\]
Some \(A_e\) is therefore nonzero with \(e\ge m+1\).
Corollary~\ref{cor:characteristic-zero} replaces \(A\) by a quotient by a
monomial ideal containing every variable square, with the same Hilbert
function. A surviving degree-\(e\) monomial is squarefree, and all its
\(2^e\) divisors survive. Consequently
\(\deg\Gamma\ge2^e\ge2^{m+1}\). This proves, in characteristic zero,
the numerical assertion \(\mathrm{IV}_m\) of \cite[Section~3]{EGH93}.

Now suppose \(X\) does not contain \(\Omega\) and \(k\le r-1\).
Put \(Z=\Omega\cap X\), and let \(\Gamma\) be residual to \(Z\) in
\(\Omega\). If \(I_Z\) and \(I_\Omega\) are their saturated homogeneous
ideals, the residual scheme is defined by
\(I_\Gamma=(I_\Omega:I_Z)\). Set \(m=r-k-1\). Residual duality for the
complete intersection of \(r\) quadrics gives
\[
 \deg\Gamma-H_\Gamma(m)
   =\dim_F(I_Z/I_\Omega)_k>0.
\]
The equality is the modern Cayley--Bacharach theorem in
\cite[Section~3, p.~195]{EGH93}; its complementary degrees sum to
\(\sum_{i=1}^r2-r-1=r-1\).
The strict inequality holds because the defining equation of \(X\)
belongs to \(I_Z\) but not to \(I_\Omega\).
This theorem may be applied over an algebraic closure: flat field
extension preserves lengths, Hilbert functions, residual schemes,
and scheme containment. Thus \(\Gamma\) fails to impose independent
degree-\(m\) conditions. The bound just proved and the residual degree
identity give
\[
 \deg(\Omega\cap X)=2^r-\deg\Gamma
 \le 2^r-2^{r-k}.
\]
For \(k=r\), the same inequality merely says that a proper closed
subscheme of \(\Omega\) has integer degree at most \(2^r-1\).
\end{proof}
